\documentclass[11pt]{article}
\usepackage[a4paper,margin=23mm]{geometry}
\usepackage[T1]{fontenc}
\usepackage{lmodern,amsmath,amssymb,amsthm,mathtools,microtype,booktabs}
\usepackage[colorlinks=true,linkcolor=blue,citecolor=blue,urlcolor=blue]{hyperref}
\newcommand{\E}{\mathbb E}
\newcommand{\Wick}{\mathcal W}
\newcommand{\tr}{\operatorname{tr}}
\newcommand{\norm}[1]{\left\lVert#1\right\rVert}
\newtheorem{theorem}{Theorem}
\newtheorem{lemma}{Lemma}
\title{One selected state for higher Wick descendants\\
\large Boundary sources, charged towers and fixed-degree energy pinching}
\author{Alan Day \and Chengheng\thanks{AI research assistant (OpenAI).}}
\date{8 November 2026\quad Local research draft}
\begin{document}\maketitle
\begin{abstract}
We extend the single-excitation boundary identity from the Hamiltonian,
size fluctuation and quadratic curvature to disjoint marked Wick blocks
of arbitrary fixed lengths, including neutral marked subblocks and
whole-composite Gaussian-chaos projections. The proof reduces every
source collision to at most four endpoints. A finite exhaustive rational
certificate verifies its local algebra. The complete charged Wick tower
has the same reduced density $1-\sqrt q+\sqrt q\,h^2$. Combining the
all-unrestricted-source identity with existing charged spectral estimates
transfers exact energy pinching to each fixed descendant in the selected
state. The degree-uniform ordinary charged bridge remains a separate task.
\end{abstract}

\section{Model and admissible expressions}
Use standard CAR, the overlapping Hermitian complex Gaussian $p$-body
Hamiltonian $H$, one common disorder, a fixed half-filled Slater vector
$O$, and $m=N/2$. Set $\lambda_N=p^2/N\to\lambda>0$,
$C_Q=\E H^2$, and let $C_m$ be its sector value. Define
\begin{align*}
 DX&=QX-\sum_i c_i^\dagger Xc_i,&
 \mathcal LX&=\tfrac12\bigl(NX-\sum_i(c_i^\dagger Xc_i+c_iXc_i^\dagger)\bigr),\\
 h_N&=H_m/\sqrt{C_m},&
 W_N&=C_m^{-1/2}(p-\mathcal L-\lambda_N)H|_m.
\end{align*}
Loss maps act before restricting charge. The normalization $C_m$ is a
fixed target-sector scalar, not the operator $C_Q$ inside a loss map.
$M_N^\dagger$ is the previously normalized same-disorder occupied-to-empty
block, with $\E\langle O|M_NM_N^\dagger|O\rangle=1$.
Time is dimensionless $h$-time and $\hbar=1$.

In an ordered list of $H$ slots, a mark $(a,b,T)$ means the deficit
$((b-a)p-T)$ on that entire interval, for $T=D,\mathcal L$.
Marked intervals are disjoint and contain no complete Gaussian pair.
Larger Wick regions can enforce stronger exclusions. A chaos region
$(a,b,d)$ keeps exactly $(b-a-d)/2$ internal pairs. These constraints
implement, for example,
\[
 \Wick_d\!\left[H^a\bigl(((vp-\mathcal L)Z_v)\diamond Z_w\bigr)\right],
 \qquad Z_v=\Wick_v(H^v).
\]
Here $\diamond$ is coupling Wick multiplication. An arbitrary fixed
scalar shift of each deficit and finite linear combinations are allowed.
Repeated overlapping loss maps are not part of this expression class.

\newpage
\section{General source and the four-endpoint reduction}
For a pairing $\pi$ with $k$ edges, the right endpoint of edge $e$ is
$C_{P_e}^\dagger C_{Q_e}$ and the left is its adjoint. A G edge is
unrestricted; an F edge has $P_e$ empty and $Q_e$ occupied relative to
$O$. Attach square-zero commuting variables $z_\alpha$ to the marks.
At a site, multiply the whole marked block before applying $1-zT_i$.
The local maps on a matrix are
\[
 D_i\begin{pmatrix}a&b\\c&d\end{pmatrix}
 =\begin{pmatrix}0&b/2\\c/2&d-a\end{pmatrix},\qquad
 \mathcal L_i\begin{pmatrix}a&b\\c&d\end{pmatrix}
 =\frac12\begin{pmatrix}a-d&b\\c&d-a\end{pmatrix}.
\]
Take the empty/occupied diagonal entry, keep the selected membership
constraint, and multiply by the Jordan--Wigner crossing sign. Denote the
result by $f_o(S;z)$, where $S$ is the set of present $P/Q$ labels.
For one label $i$, its correct-occupancy singleton is $s_i(z)$, with
$s_i(0)=1$. Set $F_o(S;z)=f_o(S;z)/\prod_{i\in S}s_i(z)$. Then
\begin{equation}
 K_\pi(z)=2\sum_{i<j}\sum_{o=0,1}F_o(\{i,j\};z)
             -2k(k-1)-2k_G .\label{K}
\end{equation}
With $\gamma=e^{-\lambda}$ and $q=\gamma^4$, a diagram contributes
\begin{equation}
 \gamma^{b_{\rm mix}}[z_1\cdots z_r]
 \exp\!\left(\lambda K_\pi(z)-\sum_\alpha\mu_\alpha z_\alpha\right),
 \qquad K_\pi(0)=-4\nu_\pi .\label{weight}
\end{equation}
$\nu_\pi$ counts crossings and nestings with an inner F edge;
$b_{\rm mix}$ counts mixed ordinary/selected edges.

The earlier fixed-Slater source derivation \cite{composite} extends
without a length-two assumption. No pair inside a mark means every
single present label gives only an odd local factor there. Thus
$\prod_i s_i^p\exp[p\sum_\alpha(b_\alpha-a_\alpha)z_\alpha]=1$ exactly.
For each fixed layout, normalized defect coefficients have bounded
nilpotent norm $B$. The old falling-factorial domination applies with
activity $B m(p/m)^{|S|}$: its tails are summable, triples disappear
after domination, and backward selected edges require a defect of
order $p$. Wick and chaos restrictions only delete pairings.
This establishes \eqref{weight} for every fixed admissible expression.

\begin{lemma}[Finite local reduction]
Every source-dependent summand of \eqref{K} is an own-edge collision
with two endpoints or a two-edge collision with four endpoints.
Deleting all other endpoints preserves its normalized local polynomial.
\end{lemma}
\begin{proof}
Only the two present labels can supply nonidentity site matrices.
An empty marked interval contributes $T_i(I)=0$. Each surviving marked
interval becomes a contiguous interval among the retained endpoints;
disjointness and absence of an internal pair survive. The singleton
denominator and crossing sign use exactly the same present labels.
An F edge touching an extreme endpoint remains extreme after deletion.
Thus the ordered local cases have no dependence on the original lengths.
\end{proof}

\newpage
\section{Local identities and the selected density}
For an own-edge collision, let $\ell_1,\ell_2$ be the source variables
at its endpoints when the corresponding mark is of type L, and zero
otherwise. The source-dependent contribution is
\[
 K_{\rm own}^{G}=\ell_1+\ell_2+\ell_1\ell_2,
 \qquad K_{\rm own}^{F}=0.
\]
The endpoints cannot belong to the same mark. In a boundary-selected
word $M B M^\dagger$, at least one endpoint of every F edge is outside
all interior marks, so its own-edge difference is just
$-\ell_1-\ell_2$.

For distinct edges, replacing a G edge that touches a selected exterior
endpoint by F leaves every nonconstant collision coefficient unchanged.
Also, an all-G L singleton placed at an exterior endpoint has no
source-containing two-edge collision. These assertions follow by
substitution of the two local matrices on the preceding page. The
exhaustive certificate enumerates the three two-edge orders (disjoint,
crossed, nested), whether the left/right global extreme is present, and
all legal contiguous marked intervals of either type. Its counts are
\begin{center}
\begin{tabular}{lr}\toprule
Local layouts (two or four endpoints)&248\\
Selected two-edge source equalities&241\\
Selected own-edge source equalities&489\\
Exterior L two-edge zero identities&241\\
Exterior L own-edge identities&489\\\bottomrule
\end{tabular}
\end{center}
All nilpotent masks are checked with exact rational arithmetic; the
complete case inventory is saved with the certificate. This is a finite
algebraic verification used together with the reduction lemma, not an
extrapolation from a finite maximum field degree.

Let $n_\alpha$ count selected-boundary attachments inside interior
L mark $\alpha$. There is no inner-F nesting, and consequently
\begin{equation}
 K_F(z)=K_G(z)-\sum_{\alpha:\,L}n_\alpha z_\alpha.\label{boundary}
\end{equation}
For the all-G word, add L singleton sources $a,b$ at the two boundaries.
The own-edge rule and two-edge zero identity give, when they are paired
to distinct interior endpoints,
\[
 K_G(z,a,b)-a-b=K_G(z)+a\ell_0+b\ell_L.
\]
Here $\ell_0,\ell_L$ are the interior L variables at the respective
partners, or zero. For a direct exterior pair the right side is instead
$K_G(z)+ab$. Insert $\chi=h-W$: each boundary applies
$(1-\partial_a)$ or $(1-\partial_b)$ at zero to the exponential.
The attached weight becomes $(1-\lambda\ell_0)(1-\lambda\ell_L)$,
exactly \eqref{boundary}. The direct exterior pair has weight
$1+\lambda$, whereas the selected direct pair has weight one.

\begin{theorem}[Fixed-degree boundary identity]
For every admissible expression $B$ and finite linear combination thereof,
with $s=\gamma^2=\sqrt q$,
\begin{equation}
 \sigma_1(B)=[1-s(1+\lambda)]\tau(B)+s\tau(\chi B\chi),
 \qquad \chi=h-W.\label{density}
\end{equation}
Here $\tau$ denotes the limiting covariant central trace and $\sigma_1$
the limiting normalized single-selected expectation. In tracial moment
notation the density is $\rho_1=s\chi^2+[1-s(1+\lambda)]I$.
\end{theorem}
The attached diagrams carry $\gamma^2$ and the direct diagrams carry
one, which proves \eqref{density}. This statement is on polynomial
moments; it does not assert boundedness of all higher generators.

\newpage
\section{The charged tower and its exact pinching}
If $B$ contains only D marks and ordinary energy factors, both
$\ell_0,\ell_L$ vanish. Therefore
\begin{equation}
 \sigma_1(B)=(1-s)\tau(B)+s\tau(hBh)=\tau(f(h)B),
 \qquad f(h)=1-s+s h^2.\label{charged-density}
\end{equation}
This includes arbitrary fixed mixed words in all charged Wick
descendants, not only the quadratic curvature. In particular the
polynomial functional is invariant under the Hamiltonian derivation.
Finite-time invariance for fixed insertions follows from the energy
tail estimates below. This is not a temperature assignment.

Let $H_d^{(q)}$ be the monic q-Hermite polynomial with variance one,
$P_d=[d]_q!$, and $m_d=q\partial_q\log P_d$. Use the finite-$N$ field
\begin{equation}
 J_{d,N}=C_m^{-d/2}(dp-D)\Wick_d(H^d)|_m
       -\lambda_N[d(d-1)-4m_d(q_N)]H_d^{(q_N)}(h_N),
 \quad q_N=e^{-4\lambda_N}.\label{Jd}
\end{equation}
For d=2 this is the preceding unnormalized curvature $j_N$.
All normalization and regression coefficients are fixed before the
double-scaling limit. The q-Hermite term commutes exactly with $h_N$.

For all-G diagrams, the full-trace local ratio is
$2\sum_o F_o$, and both constant backgrounds are $-2k^2$.
Thus the entire central marked-source exponent equals the earlier
full-trace exponent, for arbitrary legal marks. Schur covariance
equates the annealed fixed-Slater covariant expectation to normalized
sector trace. This directly transfers the existing fixed-degree
charged joint moments and spectral argument \cite{joint}; no extraction
from a binomial charge average is used.

In particular, for fixed d,e the limiting two-energy measure
\[
 \nu_{de,N}(A,B)=\E\tr_m\!\bigl[
 J_{d,N}^\dagger P_{h_N}(B)J_{e,N}P_{h_N}(A)\bigr]
\]
has the previously derived $L^2(\mu_q\otimes\mu_q)$ density.
For normalized energy polynomials, the old bound reads, with $D=d+e$,
\[
 |M_{ab}^{de}|\le K_D(a+b+1)^D
 q^{\max(0,\min(a,b)-D)}\,\mathbf1_{|a-b|\le D}.
\]
It is square summable for fixed d,e. The compact energy support is
$[-2/\sqrt{1-q},2/\sqrt{1-q}]^2$; $\mu_q$ has no atoms, so the
energy diagonal has zero mass. Central fixed higher moments give
tightness and uniform integrability. Portmanteau then yields
\begin{equation}
 \norm{\mathcal P_{h_N}J_{d,N}}_{2,m}\longrightarrow0,
 \qquad \mathcal P_hY=\sum_E P_h(\{E\})YP_h(\{E\}).\label{pinch}
\end{equation}
Higher central moments and the Schatten contraction of pinching upgrade
this to every fixed finite Schatten exponent. Complete degenerate
blocks are retained. This takes finite-$N$ exact pinching first and
does not exchange a long-time averaging limit with double scaling.

\newpage
\section{Selected actual observables and the remaining time interface}
The previously proved fixed-polynomial preparation estimate
\cite{polynomial} gives, for any fixed preparation $P_N$ with bounded
$\E\norm{P_NO}^2$,
\[
 \E\norm{X_NP_NO}^2\le C_P\norm{X_N}_{4,m}^2.
\]
Apply it to $X_N=\mathcal P_{h_N}J_{d,N}$. Since the new source theorem
also bounds every fixed preparation norm, \eqref{pinch} implies
\[
 \E\norm{\bigl[(1-\mathcal P_{h_N})J_{d,N}-J_{d,N}\bigr]P_NO}^2
 \longrightarrow0.
\]
For single-selected covariant blocks, polarization, trace Holder and
fixed higher moments similarly transfer any fixed finite number of
these replacements. Fixed real-time energy phases are included by
first approximating them in central trace norms and then applying the
preparation bound. Indeed Gaussian energy moments obey
$\norm{h_N^n}_{r,m}\le K_r^n\sqrt{n!}$, making their Taylor tails
summable at every fixed compact time interval after Holder.

Writing $a_{d,N}=(1-\mathcal P_{h_N})J_{d,N}$, the charged two-point
matrix in the single-selected state is therefore
\begin{equation}
 \lim_N\E\langle O|M_N a_{d,N}(t)a_{e,N}M_N^\dagger|O\rangle
 =\int f(E)e^{it(E-E')}\,d\nu_{de}(E,E').\label{kernel}
\end{equation}
The input/column energy is E in this convention. Its positive matrix
measure is $f(E)\nu_{de}$, and the D-only state is stationary. We do
not replace every descendant's marginal by the tracial marginal:
that further property was proved specifically for the curvature field.

The previous tagged-site proof \cite{tails} already supplies, sector
by sector, for fixed r and a fixed number k of deficits,
\[
 \norm{\Delta_{T_1}\cdots\Delta_{T_k}Z_n}_{r,m}
 \le K_{k,r}(n+1)^{2k}L_{k,r}^n\sqrt{n!}
\]
in central charge ranges. It uses restricted covariance row/column
bounds and keeps CAR signs until matrix contraction. Coupling OU
projections, normalized covariance CP chains and their telescoping
estimates also operate sectorwise. These estimates support pure/local
representatives, including neutral Wick subblocks.

The ordinary charged coefficient needed by the actual full bilocal is
\[
 G_{n,N}=[p\mathcal N-D-\lambda_N\mathcal N(\mathcal N-1)]H^n.
\]
It permits internal Gaussian pairs in its marked interval, outside
the present theorem. Its existing full-Fock tail proof uses the
Clifford Riesz estimate \cite{belp}, whose trace is not the fixed-charge
trace. Restriction to $m=N/2$ is not automatic. A central strong bridge
and a degree-uniform actual tail, or an equivalent two-cutoff argument,
are still needed to identify the full selected bilocal. Fixed-degree
preparation constants cannot be summed over growing degrees without
additional control. The previously completed full-trace bilocal four-point
result is unchanged.

\newpage
\section{Evidence, reuse and continuation}
The new pure module \texttt{slater\_higher\_wick\_sources.py} accepts
slot lengths, disjoint marks, Wick exclusion regions, chaos regions,
scalar centers and trace/selected/$\chi$ modes. It is built on the
already checked local CAR engine. Its coefficient exponential retains
all nilpotent source masks. The fresh checker saves the complete finite
local case inventory, 30 higher-layout density comparisons at two
independent rational $(\gamma,\lambda)$ points, and 18 comparisons
against the previous H/W/J evaluator. Higher examples reach eight
internal slots; they include two long D marks, neutral L subblocks,
several simultaneous marks, and projected whole composites.

The rational points check algebraic identities; positivity and the
physical state interpretation use $\gamma=e^{-\lambda}$. Finite examples
alone do not prove all degrees: the finite local certificate is made
exhaustive by the reduction lemma. The convergence and pinching claims
also use the specified earlier estimates, not just this checker.
Input hashes, Python version and exact arithmetic are saved with the
new output. Successful prior scientific drivers and manifests are
left unchanged.

The source extension is the new construction in this draft. Fixed-charge
complex-SYK chord methods \cite{bnr}, the previous local source theorem,
the charged spectral bound and the preparation-density estimate are
reused with their original scopes. The unresolved ordinary central
bridge is now the direct next step. Bulk spatial and mass generators,
and finite-G dynamics, remain beyond this model-level result.

\begin{thebibliography}{9}
\bibitem{bnr} M. Berkooz, V. Narovlansky and H. Raj,
\emph{Complex Sachdev--Ye--Kitaev model in the double scaling limit},
\href{https://arxiv.org/abs/2006.13983}{arXiv:2006.13983v2},
especially Sections 3 and 5.
\bibitem{belp} L. Ben Efraim and F. Lust-Piquard,
\emph{Poincare type inequalities on the discrete cube and in the CAR algebra},
\href{https://arxiv.org/abs/math/0702233}{arXiv:math/0702233v1},
Sections 6.1 and 6.4, Theorem 6.8.
\bibitem{composite} Alan Day and Chengheng,
\emph{Same-J Slater composite sources} and
\emph{A selected excitation and its curvature dynamics},
local project drafts, 8 November 2026; companion notes
\texttt{susskind-slater-composite} and \texttt{susskind-slater-selected-state}.
\bibitem{joint} Alan Day and Chengheng,
\emph{Joint memory regression}, local project proof,
7 November 2026, Sections 5--7;
\texttt{susskind-memory-joint-regression-proof.md}.
\bibitem{polynomial} Alan Day and Chengheng,
\emph{Fixed polynomial Slater preparations}, local project draft,
8 November 2026; \texttt{susskind-slater-polynomial-density}.
\bibitem{tails} Alan Day and Chengheng,
\emph{Bilocal fourth-tail diagram and operator drafts}, local project proofs,
7 November 2026; pure marked-block estimates and sectorwise CP bounds.
\end{thebibliography}
\end{document}
